\documentclass[10pt,a4paper]{article}
\usepackage[margin=12mm]{geometry}
\usepackage[T1]{fontenc}
\usepackage[utf8]{inputenc}
\usepackage{lmodern}
\usepackage{amsmath,amssymb,booktabs,tabularx,array,xcolor,tcolorbox,microtype}
\usepackage{enumitem}
\pagestyle{empty}
\setlength{\parindent}{0pt}
\setlength{\parskip}{3pt}
\definecolor{navy}{RGB}{18,48,82}
\definecolor{green}{RGB}{29,112,76}
\definecolor{amber}{RGB}{171,111,0}
\definecolor{red}{RGB}{151,45,45}
\definecolor{softblue}{RGB}{239,245,251}
\definecolor{softgray}{RGB}{246,247,249}
\newtcolorbox{proofbox}[1][]{colback=softblue,colframe=navy,boxrule=.7pt,arc=1.5mm,left=2mm,right=2mm,top=1.5mm,bottom=1.5mm,#1}
\newtcolorbox{statusbox}[1][]{colback=softgray,colframe=black!45,boxrule=.6pt,arc=1.5mm,left=2mm,right=2mm,top=1.5mm,bottom=1.5mm,#1}
\newcommand{\PASS}{\textcolor{green}{\textbf{PASS}}}
\newcommand{\OPEN}{\textcolor{amber}{\textbf{OFFEN}}}
\newcommand{\FAIL}{\textcolor{red}{\textbf{FAIL}}}
\begin{document}

{\Large\bfseries EABC -- Beweisfahrplan für das Kolloquium}\hfill{\small Handout v8}\\[-1mm]
{\small Arithmetik, ternäre Gates und platonisch-quaternionische Geometrie -- Status strikt getrennt}
\vspace{2mm}
\hrule
\vspace{3mm}

\begin{proofbox}
\textbf{Leitprinzip.} Keine Struktur wird durch Analogie aufgewertet. Jeder Übergang wird als eigener, falsifizierbarer Pfeil behandelt:
\[
\text{arithmetische Invariante}
\;\stackrel{?}{\longrightarrow}\;
\text{Grad-/Faserstruktur}
\;\stackrel{?}{\longrightarrow}\;
\text{Fano-/Gitterstruktur}.
\]
\end{proofbox}

\textbf{1. Quadratische Reziprozität: exakt lokalisierte Reichweite}

Für Primzahlen $p>3$ wird die EABC-Klasse durch
\[
(\mathbb Z/12\mathbb Z)^\times=\{1,5,7,11\}\cong V_4
\]
organisiert. Der reziproke Symmetriedefekt reduziert sich exakt auf das mod-$4$-Bit
\[
c(p,q)=b(p)b(q)\in\mathbb F_2.
\]
Damit entsteht ein sauberer binärer Superbraiding-Sektor. Der daraus gebildete ternäre Orientierungsdefekt
\[
T(x,y,z)=b_xb_y+b_yb_z+b_zb_x
\]
ist jedoch paarweise rekonstruierbar und verletzt die Alternativität, denn $T(x,x,z)=b_x$. Daher:
\[
\boxed{\text{quadratische Reziprozität}\not\Rightarrow\text{oktonionischer Assoziator}.}
\]

\textbf{2. Gate $G_5$: genuine ternäre Arithmetik}

\begin{center}
\small
\begin{tabularx}{\textwidth}{>{\bfseries}l X l}
\toprule
Gate & Prüfgegenstand & Status\\
\midrule
$G_{5.0}$ & paarweise Legendre-Trivialität / Eingang in den ternären Testsektor & exakt prüfbar\\
$G_{5.1}$ & Rédei-Symbol als unabhängige ternäre arithmetische Invariante & Referenzanker vorhanden\\
$G_{5.2}$ & Abstieg auf EABC-Grade & \OPEN\\
$G_{5.3}$ & volle ternäre/Fano-kompatible Struktur & \OPEN\\
$G_{5.4}$ & exakter Match mit v26-Fano-Assoziator & \OPEN\\
\bottomrule
\end{tabularx}
\end{center}

\begin{statusbox}
\textbf{Referenzpaar.} $(13,17,97)$ scheitert bereits an $G_{5.0}$. Das Tripel $(5,41,61)$ besitzt paarweise triviale Legendre-Daten und dient mit dem klassischen Rédei-Wert $[5,41,61]=-1$ als ternärer Anker. Daraus folgt \emph{nicht} bereits ein Fano-Descent.
\end{statusbox}

\textbf{3. Klassische geometrische Kette}
\[
\boxed{
\text{Ikosaeder}\to A_5\to 2I\subset\mathbb H\to H_4\to\text{Ikosianen}\to E_8
}
\]
Diese Kette ist etablierte Mathematik. Offen ist nicht ihre Existenz, sondern die kanonische Kopplung der EABC-Faktorwort-/mod-$12$-Gradierung an diese Geometrie.

\vfill
{\footnotesize Statussprache: \PASS\ = mathematisch/auditiert gesichert; \OPEN\ = explizite Brücke fehlt; \FAIL\ = vorgeschlagene direkte Identifikation falsifiziert.}

\newpage

{\Large\bfseries Seite 2 -- Geometrische Reduktion an $G_{\mathrm{ID}}$}\hfill{\small Maschinenexakter Gruppentheorie-Audit}
\vspace{2mm}
\hrule
\vspace{3mm}

\begin{proofbox}
\textbf{Ausgangspunkt.} Eine Kleinsche Vierergruppe $V_4\le A_5$ wirkt frei auf den 12 Ikosaederecken. Daher zerfällt die Eckenmenge in genau drei freie $V_4$-Orbits zu je vier Punkten. Eine $V_4$-äquivariante EABC-Beschriftung ist auf jedem Orbit durch einen Shift $s_i\in V_4$ bestimmt.
\end{proofbox}

\begin{center}
\vspace{1mm}
{\Large\bfseries
$\boxed{64}$
$\;\xrightarrow{\ /V_4\ \text{Gauge}\ }\;$
$\boxed{16}$
$\;\xrightarrow{\ N_{A_5}(V_4)/V_4\cong\mathbb Z_3\ }\;$
$\boxed{4}$}
\vspace{2mm}

\small
\begin{tabular}{c@{\qquad}c@{\qquad}c}
$4^3$ äquivariante Labelings & Translationsklassen & Normalisator-Fixklassen\\
\end{tabular}
\end{center}

\vspace{3mm}
\textbf{Exakter Audit.}
\[
|A_5|=60,\qquad |V_4|=4,\qquad |N_{A_5}(V_4)|=12.
\]
Auf den 16 Gauge-Klassen wirkt der Quotient
\[
N_{A_5}(V_4)/V_4\cong C_3.
\]
Die Orbitzerlegung lautet exakt
\[
\boxed{1+1+1+1+3+3+3+3=16.}
\]
Äquivalent auf den 64 ursprünglichen Beschriftungen:
\[
\boxed{4+4+4+4+12+12+12+12=64.}
\]

\begin{center}
\begin{tcolorbox}[width=.88\textwidth,colback=softblue,colframe=navy,boxrule=.8pt,arc=1.5mm]
\centering
{\large\bfseries Die Fixpunkt-Triade}\\[1mm]
\[
\boxed{4\text{ normalisator-fixe Gauge-Klassen}}
=\boxed{1\text{ trivial}}\;\sqcup\;\boxed{3\text{ nichttriviale Kandidaten }\mathcal C_1,\mathcal C_2,\mathcal C_3}.
\]
\end{tcolorbox}
\end{center}

\textbf{Gate-Status nach dem Audit}

\begin{center}
\small
\begin{tabularx}{\textwidth}{>{\bfseries}l X l}
\toprule
Gate & Aussage & Status\\
\midrule
$G_{\mathrm{ID},4a}$ & $V_4$-äquivariante $3+3+3+3$-Beschriftungen existieren & \PASS\\
$G_{\mathrm{ID},4b}^{\rm geom}$ & Geometrie reduziert $16$ Gauge-Klassen auf $4$ Fixklassen & \PASS\\
$G_{\mathrm{ID},4c}$ & arithmetischer Selektor wählt kanonisch unter $\mathcal C_1,\mathcal C_2,\mathcal C_3$ & \OPEN\\
$G_{\mathrm{ID},5}$ & kompatibler Faktorwort-$\to$Quaternionen-/Ikosian-/E8-Lift & \OPEN\\
\bottomrule
\end{tabularx}
\end{center}

\begin{statusbox}
\textbf{Verbleibende Forschungsfrage.} Die Geometrie liefert drei nichttriviale, normalisator-fixe Kandidaten. Zu prüfen ist, ob eine unabhängige arithmetische Invariante -- etwa Rédei-/höhere ternäre Daten -- eine wohldefinierte, repräsentantenunabhängige Auswahl
\[
\{\mathcal C_1,\mathcal C_2,\mathcal C_3\}\stackrel{?}{\longrightarrow}\mathcal C_{\rm arith}
\]
erzeugt. Ein solcher Selektor wäre ein neues Theorem; er wird nicht vorausgesetzt.
\end{statusbox}

\begin{proofbox}
\textbf{Prüfungssatz für die Kommission.} Die Reduktion $64\to16\to4$ ist ein endliches, maschinenexakt reproduziertes Gruppentheorie-Resultat. Die Auswahl einer der drei nichttrivialen Fixklassen und ihr möglicher $E_8$-Lift bleiben ausdrücklich getrennte offene Gates.
\end{proofbox}

\vfill
{\footnotesize Reproduzierbarer Audit: $A_5/C_5$ als 12-Eckenmenge, $V_4$-Normalisator in $A_5$, kombinierte Normalisator-/Gauge-Wirkung. Keine Aussage über thermische oder dynamische Stabilität wird aus dem Stabilisator allein abgeleitet.}


\newpage

{\Large\bfseries Seite 3 -- Additive Conservation Sector}\hfill{\small Additive Bilanz und Faktorwort strikt getrennt}
\vspace{2mm}
\hrule
\vspace{3mm}
\small

\begin{proofbox}
\textbf{Definition (additive Ladung).} Für einen endlichen Zustand $X=\{x_1,\dots,x_k\}$ aus zugelassenen arithmetischen Labels definieren wir
\[
Q(X):=\sum_{j=1}^{k}x_j.
\]
Ein Übergang $X\to Y$ heißt \emph{additiv zulässig}, wenn
\[
\boxed{Q(X)=Q(Y).}
\]
Dies ist zunächst eine definierte kombinatorische Invariante. Eine Interpretation als Energie- oder Noether-Ladung erfordert zusätzlich eine Dynamik, eine Wirkung und eine zugehörige kontinuierliche Symmetrie und wird hier nicht vorausgesetzt.
\end{proofbox}

\textbf{4. Additiver und multiplikativer Sektor}

Die additive Bilanz wird vom bisherigen multiplikativen Faktorwortsektor getrennt geführt:
\[
\boxed{Q(X)=\sum_{x\in X}x}
\qquad\text{gegenüber}\qquad
\boxed{W(N)=(p_1,\ldots,p_{\Omega(N)})}.
\]
Im Faktorwortsektor bleiben ausschließlich die atomaren Primfaktoren mit Multiplizität erhalten; zusammengesetzte Teilprodukte erzeugen keine neuen Generatoren. Der additive Sektor klassifiziert dagegen zulässige Umordnungen beziehungsweise Zerlegungen mit konstantem $Q$.

\begin{center}
\small
\begin{tabularx}{\textwidth}{>{\bfseries}p{3.2cm} X X}
\toprule
Objekt & Mathematische Rolle & Status im Modell\\
\midrule
$Q(X)$ & additive Summeninvariante & definitorisch exakt\\
$W(N)$ & eindeutiges multiplikatives Faktorwort & arithmetische Grundstruktur\\
$1$ & multiplikatives neutrales Element; keine Primzahl & exakt\\
$1$ im additiven Sektor & optionaler Bilanzausgleich in zugelassenen Zerlegungen & Modellkonvention\\
``Energie'' von $Q$ & physikalische Interpretation & nicht vorausgesetzt\\
\bottomrule
\end{tabularx}
\end{center}

\textbf{5. Die Doppelrolle der $1$}

Die Zahl $1$ wird ausdrücklich \emph{nicht} als Primzahl behandelt. Multiplikativ gilt
\[
1\cdot p=p,
\]
und für das leere Faktorwort kann konsistent
\[
W(1)=\varnothing
\]
gesetzt werden. Additiv darf $1$, sofern die gewählte Zustandsklasse dies zulässt, als kleinster Bilanzausgleich auftreten, etwa
\[
2+2=3+1.
\]
Diese Gleichung ist eine additive Partition von $4$; sie beweist weder eine Primzahlerzeugung noch eine physikalische Selbstenergie.

\begin{statusbox}
\textbf{Epistemische Grenze.} Die Formulierung
\[
X\to Y\;\Longrightarrow\;Q(X)=Q(Y)
\]
ist ein mathematisch testbarer Erhaltungskanal innerhalb des definierten Übergangssystems. Erst ein separat konstruierter Hamiltonoperator $H$, eine Wirkung $S$ oder eine kontinuierliche Symmetrie könnte eine Noether- beziehungsweise Energieinterpretation rechtfertigen.
\end{statusbox}

\textbf{6. Primzahlzwillinge als Observable, nicht als Postulat}

Für ungerade Primzahlen ist $\Delta p=2$ der kleinste mögliche positive Abstand zweier verschiedener Primzahlen. Das Modell darf daher untersuchen, ob seine \emph{eigenständig definierte} Dynamik Übergänge mit
\[
\Delta p=2
\]
statistisch bevorzugt. Die zulässige Forschungsfrage lautet
\[
\boxed{\text{Erzeugt die arithmetische Übergangsdynamik einen Bias bei }\Delta p=2\text{?}}
\]
Eine positive numerische Antwort wäre zunächst ein statistischer Befund des Modells und kein Beweis der Primzahlzwillingsvermutung.

\begin{proofbox}
\textbf{Prüfungssatz für die Kommission.} Der Additive Conservation Sector ergänzt die Faktorwortarchitektur um eine exakt definierte Summeninvariante, ohne die $1$ zur Primzahl oder zu einer postulierten Vakuumenergie aufzuwerten. Additive Bilanz, multiplikative Faktorisierung und physikalische Interpretation bleiben drei getrennte Ebenen.
\end{proofbox}

\vfill
{\footnotesize Methodischer Status: mathematische Definitionen und arithmetische Identitäten = exakt; dynamischer $\Delta p=2$-Bias = testbare Hypothese; Noether-/QFT-Deutung = nur nach separater Konstruktion einer geeigneten Dynamik.}


\newpage

{\Large\bfseries Seite 4 -- Airy-Kalibrierung des lokalen Cusp-Operators}\hfill{\small Exakte Referenzlösung}
\vspace{2mm}
\hrule
\vspace{3mm}
\small

\begin{proofbox}
\textbf{Lokales Referenzproblem.} Für
\[
H_a=-\frac{d^2}{dT^2}+a|T|,\qquad a>0,
\]
führt die Skalierung $T=a^{-1/3}x$ exakt auf
\[
H_a=a^{2/3}\left(-\frac{d^2}{dx^2}+|x|\right).
\]
Damit gilt für alle Eigenwerte
\[
\boxed{E_n(a)=a^{2/3}E_n(1).}
\]
Für $a=8$ ist die universelle Skaleneinheit daher $8^{2/3}=4$.
\end{proofbox}

\textbf{7. Parität und Airy-Nullstellen}

Auf der ganzen Geraden zerfällt das Spektrum nach Parität. Für $T>0$ erfüllt die Eigenfunktion die Airy-Gleichung. Gerade Zustände gehorchen $\psi'(0)=0$ und werden durch Nullstellen von $\operatorname{Ai}'$ bestimmt; ungerade Zustände gehorchen $\psi(0)=0$ und werden durch Nullstellen von $\operatorname{Ai}$ bestimmt. Mit den positiven Beträgen der negativen Nullstellen $|a_k'|$ bzw. $|a_k|$ gilt insbesondere
\[
E_0=4|a_1'|,\qquad E_1=4|a_1|.
\]
Numerisch
\[
|a_1'|\approx1.018793,\qquad |a_1|\approx2.338107,
\]
also
\[
\boxed{E_0\approx4.07517},\qquad
\boxed{E_1\approx9.35243}.
\]
Die Abweichung von $E_0$ gegenüber der bloßen Skaleneinheit $4$ ist daher kein Gitterfehler, sondern der exakte dimensionslose Airy-Faktor $|a_1'|$.

\begin{center}
\small
\begin{tabularx}{\textwidth}{>{\bfseries}p{2.4cm} X X}
\toprule
Größe & Analytische Referenz & Bedeutung\\
\midrule
$a^{2/3}$ & $4$ bei $a=8$ & universelle Skaleneinheit\\
$E_0$ & $4|a_1'|\approx4.07517$ & gerader Grundzustand\\
$E_1$ & $4|a_1|\approx9.35243$ & erster ungerader Zustand\\
$\Delta E$ & $4(|a_1|-|a_1'|)\approx5.2773$ & fundamentale Spektrallücke des Cusp-Modells\\
\bottomrule
\end{tabularx}
\end{center}

\textbf{8. Exakte lokale Spektrallücke}

Die erste Lücke ist analytisch bestimmt durch
\[
\boxed{\Delta E=E_1-E_0=4\bigl(|a_1|-|a_1'|\bigr)\approx5.2773.}
\]
Sie ist damit keine Fitgröße des numerischen Solvers. Eine numerische Rechnung am diskretisierten Operator dient hier als Reproduktion einer geschlossenen Referenzlösung.

\begin{statusbox}
\textbf{Cusp statt Singularität.} Das Potential $8|T|$ ist endlich und stetig auf ganz $\mathbb R$, aber bei $T=0$ nicht differenzierbar. Der Ursprung ist daher ein \emph{Cusp} beziehungsweise Knick, keine divergierende Potentialsingularität.
\end{statusbox}

\textbf{9. Grenze zum vollständigen exponentiellen Confinement}

Für das im Confinement-Sektor betrachtete Potential
\[
U_8(T)=1+e^{8|T|}
\]
gilt lokal
\[
U_8(T)=2+8|T|+32T^2+\frac{256}{3}|T|^3+O(T^4).
\]
Der Operator mit $8|T|$ ist daher die führende nichtkonstante lokale Cusp-Referenz. Seine Eigenwerte sind \emph{nicht} mit dem vollständigen Spektrum von $-d^2/dT^2+U_8(T)$ gleichzusetzen.

\begin{proofbox}
\textbf{Prüfungssatz für die Kommission.} Die Numerik wird an einem analytisch exakt lösbaren Airy-Problem kalibriert. Für $a=8$ ist $a^{2/3}=4$, während die exakten Eigenwerte durch die Nullstellen von $\operatorname{Ai}$ und $\operatorname{Ai}'$ festgelegt werden; insbesondere $\Delta E\approx5.2773$. Diese Übereinstimmung validiert den lokalen numerischen Solver und die Cusp-Kalibrierung, nicht bereits das globale exponentielle Confinement oder eine weitergehende physikalische Interpretation.
\end{proofbox}

\vfill
{\footnotesize Status: Airy-Spektrum des linearen Cusp-Operators = analytisch exakt; numerische Übereinstimmung = Solver-Kalibrierung; Übertragung auf $U_8(T)=1+e^{8|T|}$ = lokale Referenz, keine Spektralidentität.}

\end{document}
